\PassOptionsToPackage{numbers, compress}{natbib}
\documentclass{article}

% NeurIPS 2026 Workshop Style - Double-Blind Workshop Track
\usepackage[dblblindworkshop]{neurips_2026}
\workshoptitle{New in ML Workshop}

\usepackage[utf8]{inputenc}
\usepackage[T1]{fontenc}
\usepackage{microtype}
\usepackage{amsmath,amssymb}
\usepackage{graphicx}
\usepackage{booktabs}
\usepackage{xcolor}
\usepackage{enumitem}
\setlist{itemsep=2pt,topsep=2pt}
\usepackage[font=small,labelfont=bf]{caption}
\usepackage{subcaption}
\captionsetup{aboveskip=4pt,belowskip=2pt}
\usepackage{tikz}
\usetikzlibrary{arrows.meta,positioning,calc}
\usepackage{pgfplots}
\pgfplotsset{compat=1.18}
\usepackage{hyperref}
\hypersetup{colorlinks=true,linkcolor=blue!70!black,urlcolor=blue!70!black,citecolor=green!50!black}
\usepackage{cleveref}

% Float placement — NeurIPS template defaults are used

\title{Anomaly-Driven Correction Discovery: Constraining Hypothesis Space Beats Scaling Compute in Noisy Symbolic Regression}

% Author block is auto-rendered anonymously by neurips_2026.sty [dblblindworkshop]
\author{Anonymous Author(s)}

\begin{document}
\maketitle

\begin{abstract}
Symbolic regression (SR) frameworks such as PySR, AI~Feynman, and PhySO discover physical laws from scratch (\emph{tabula rasa}). Under realistic observational noise, however, their structural recovery degrades rapidly and \emph{non-monotonically}: doubling PySR's candidate budget does not reliably improve functional class recovery. We argue that scientific discovery is rarely \emph{tabula rasa}; it is \emph{correction-first}. Given a known classical baseline $f_{\text{base}}(X)$ and anomalous residuals, the physical objective is discovering the minimal valid correction $\Delta(X)$. We introduce \textbf{Anomaly-Driven Correction Discovery (ADCD)}, which restricts search to a physics-anchored correction subspace $\mathcal{D} \subset \mathcal{F}$ and cascades three analytical gates (AST complexity, Buckingham-$\Pi$ dimensional homogeneity, and asymptotic decoupling) \emph{before} JAX-gradient optimization and Bayesian Information Criterion reranking. Across 16 random seeds on a 9-scenario synthetic benchmark at $5\%$ noise, ADCD achieves a mean structural recovery rate of $77.1\%$ (range: $69.4\%$--$94.4\%$) versus $11.1\%$ for PySR under matched computational budget and comparable wall-clock time, establishing that the advantage lies in accuracy rather than speed. On four real physical scenarios, ADCD recovers the correct functional class in all cases. Applied to the SPARC dataset ($3{,}342$ observations across $171$ disc galaxies), ADCD autonomously discovers a member of the Simple-MOND interpolating family, achieving cross-validated NMSE of $0.386$ versus $0.673$ for the zero-parameter Simple MOND baseline.
\end{abstract}

%==================================================================
\section{Introduction}
\label{sec:intro}

When a classical physical law fails under extreme conditions, the precise scientific question is: ``what is the \emph{minimal correction} to the established law?'' The dominant SR paradigm instead answers a different, harder question: discovering $f(X) \approx y$ from a blank slate. Pioneered by Eureqa~\cite{schmidt2009distilling}, AI~Feynman~\cite{udrescu2020ai}, and PySR~\cite{cranmer2020pysr}, and extended by deep RL approaches such as DSR~\cite{petersen2021deep} and unit-constrained PhySO~\cite{tenachi2023physo}, these systems perform well in low-noise synthetic regimes. Under realistic noise, however, structural recovery collapses. Critically, this collapse is \emph{non-monotonic}: in several scenarios, PySR recovers more correct classes at $10\%$ noise than at $5\%$, indicating that its search is not calibrated to noise level.

This fragility stems from a mismatch in problem formulation. Scientific progress rarely starts from a blank slate; it \emph{corrects} existing models. Relativistic kinetic energy appends $+\frac{3}{4}mv^4/c^2$ to the Newtonian baseline; Debye screening multiplies the Coulomb potential by $e^{-r/\lambda_D}$; MOND modifies Newtonian gravity via a dimensionless interpolating function $\nu(g_{\text{bar}}/a_0)$. In each case, a trusted classical baseline is known, and the discovery task is identifying the minimal physically valid correction $\Delta$.

We present \textbf{Anomaly-Driven Correction Discovery (ADCD)}, structured around three contributions:
\begin{enumerate}[leftmargin=*]
    \item \textbf{Correction-First Formulation}: ADCD searches for a dimensionless correction $\Delta(X;\theta)$ anchored to a known baseline $y_{\text{classical}}$, reducing the effective hypothesis space from $\mathcal{F}$ to a physical correction subspace $\mathcal{D} \subset \mathcal{F}$.
    \item \textbf{Cascaded Physics Screening}: Three analytical gates (AST token complexity, Buckingham-$\Pi$ dimensional homogeneity, and asymptotic decoupling with $\Delta \to 0$) screen candidates \emph{before} any numerical fitting, eliminating structurally invalid expressions at zero optimization cost.
    \item \textbf{Multi-Seed Empirical Validation}: Evaluated across 16 random seeds and four noise levels, ADCD achieves a mean recovery rate of $77.1\%$ at $5\%$ noise versus $11.1\%$ for PySR at matched computational budget. Unlike PySR, ADCD's accuracy declines \emph{monotonically} with noise.
\end{enumerate}

%==================================================================
\section{Related Work}
\label{sec:related}

SR for physical sciences spans \emph{tabula rasa} equation discovery and domain-constrained optimization. Information-theoretic approaches include MDL~\cite{rissanen1978mdl} and SINDy~\cite{brunton2016sindy}. Neural architectures embed physical priors through PINNs~\cite{raissi2019physics}, HNNs~\cite{greydanus2019hamiltonian}, and LNNs~\cite{cranmer2020lagrangian}. Grey-box residual learning has been applied to turbulence modeling~\cite{weatheritt2016turbulence} and Universal Differential Equations~\cite{rackauckas2020universal}. SRBench~\cite{lacava2021srbench} provides standardized benchmarking. \Cref{tab:related_work} positions ADCD within this landscape.

\begin{table}[h!]
\centering
\caption{Comparison of ADCD against established symbolic regression frameworks.}
\label{tab:related_work}
\resizebox{\textwidth}{!}{
\begin{tabular}{lcccc}
\toprule
\textbf{Framework} & \textbf{Problem Formulation} & \textbf{Physical Screening} & \textbf{Compute Requirement} & \textbf{Noise Monotonicity} \\
\midrule
AI Feynman~\cite{udrescu2020ai} & Tabula rasa ($f(X) \approx y$) & Dimensional / Symmetry & Multi-core CPU / GPU & Non-monotonic \\
PySR~\cite{cranmer2020pysr} & Tabula rasa ($f(X) \approx y$) & User-defined constraints & Julia multi-core CPU & Non-monotonic \\
PhySO~\cite{tenachi2023physo} & Tabula rasa ($f(X) \approx y$) & Dimensional Analysis & Deep RL (GPU required) & Moderate \\
\textbf{ADCD (Ours)} & \textbf{Correction-first ($f_{\text{base}} + \Delta$)} & \textbf{AST + Dim + ARC Limits} & \textbf{Single-core CPU, no GPU} & \textbf{Monotonic} \\
\bottomrule
\end{tabular}
}
\end{table}

%==================================================================
\section{The ADCD Framework}
\label{sec:method}

ADCD operates as a closed-loop, physics-gated pipeline. \Cref{fig:architecture} illustrates the five-stage architecture.

\begin{figure}[h!]
\centering
\resizebox{\textwidth}{!}{
\begin{tikzpicture}[node distance=0.8cm and 0.8cm, auto, >=Latex]
    \definecolor{navydark}{RGB}{8, 60, 125}
    \definecolor{navybg}{RGB}{240, 246, 255}
    \definecolor{slateborder}{RGB}{100, 116, 139}
    \tikzset{
        box/.style={rectangle, draw=slateborder, fill=white, rounded corners=5pt, text width=4.2cm, minimum height=1.9cm, align=center, font=\small, thick, inner sep=6pt},
        accentbox/.style={rectangle, draw=navydark, fill=navybg, rounded corners=5pt, text width=4.2cm, minimum height=1.9cm, align=center, font=\small, very thick, inner sep=6pt},
        outbox/.style={rectangle, draw=navydark, fill=white, rounded corners=5pt, text width=4.2cm, minimum height=1.9cm, align=center, font=\small, very thick, inner sep=6pt},
        arrow/.style={->, >=stealth, thick, color=navydark}
    }
    \node [box] (s0) {
        \textbf{Stage 0: Input \& Residual}\\[0.3em]
        \footnotesize Classical Law: $f_{\text{base}}(X)$\\[0.15em]
        \footnotesize Data: $(X, y_{\text{obs}})$\\[0.15em]
        \footnotesize Residual: $\delta = y_{\text{obs}} - f_{\text{base}}$
    };
    \node [box, right=of s0] (s0b) {
        \textbf{Stage 0b: Proposer}\\[0.3em]
        \footnotesize \textbullet~Template Sampling\\[0.15em]
        \footnotesize \textbullet~LLM / Grammar Vocabulary\\[0.15em]
        \footnotesize \textbullet~Feature Prior Bias
    };
    \node [accentbox, right=of s0b] (s1) {
        \textcolor{navydark}{\textbf{Stage 1: Physics Screening}}\\[0.3em]
        \footnotesize \textbullet~AST Complexity Depth $\le 7$\\[0.15em]
        \footnotesize \textbullet~Buckingham-$\Pi$ Units\\[0.15em]
        \footnotesize \textbullet~ARC Limits ($\Delta \to 0$)
    };
    \node [box, below=1.2cm of s1] (s2) {
        \textbf{Stage 2: JAX Optimizer}\\[0.3em]
        \footnotesize \textbullet~JIT-Traced L-BFGS-B\\[0.15em]
        \footnotesize \textbullet~Exact AD Gradients\\[0.15em]
        \footnotesize Opt Params: $\theta=\{\theta_0, \theta_1\}$
    };
    \node [box, left=of s2] (s3) {
        \textbf{Stage 3: Occam's Selection}\\[0.3em]
        \footnotesize \textbullet~Schwarz BIC Reranking\\[0.15em]
        \footnotesize $\text{BIC} = N\ln(\text{NMSE}) + k\ln N$\\[0.15em]
        \footnotesize Complexity Penalty
    };
    \node [outbox, left=of s3] (sout) {
        \textcolor{navydark}{\textbf{Validated Output}}\\[0.3em]
        \footnotesize \textbullet~Dimensionless Correction $\Delta(X; \theta)$\\[0.15em]
        \footnotesize \textbullet~$y_{\text{true}} = f_{\text{base}} (1 + \Delta)$\\[0.15em]
        \footnotesize \textbf{Ranked Candidate}
    };
    \draw [arrow] (s0) -- (s0b);
    \draw [arrow] (s0b) -- (s1);
    \draw [arrow] (s1) -- (s2);
    \draw [arrow] (s2) -- (s3);
    \draw [arrow] (s3) -- (sout);
    \draw [arrow, dashed, color=navydark!80] (s3.north) -- (s0b.south)
        node [midway, fill=white, font=\scriptsize\itshape, inner sep=2pt] {Feedback Loop (Stuck Count)};
\end{tikzpicture}
}
\caption{System architecture of the Anomaly-Driven Correction Discovery (ADCD) framework.}
\label{fig:architecture}
\end{figure}

\subsection{Residual Extraction and Physical Gates}
Given data $(X, y_{\text{obs}})$ and classical baseline $y_{\text{classical}} = g(X; \phi)$, ADCD computes the physical residual $\delta = y_{\text{obs}} - y_{\text{classical}}$ (additive) or $\delta = y_{\text{obs}}/y_{\text{classical}} - 1$ (multiplicative). Proposed candidates are then screened through three analytical gates in sequence:
\begin{enumerate}[leftmargin=*]
    \item \textbf{AST Complexity Gate}: Candidates are parsed into Abstract Syntax Trees. Expressions exceeding token depth $7$ or total token count $20$ are rejected before evaluation.
    \item \textbf{Dimensional Guardrail (Buckingham-$\Pi$)}: Enforces unit consistency and ensures all non-linear function arguments ($\exp, \ln, \sqrt{\cdot}$) receive dimensionless ratios.
    \item \textbf{Asymptotic Limits Gate (ARC)}: Verifies $\lim_{x \to x_0} \Delta(x) = 0$ at classical boundary conditions, ensuring physical decoupling in regimes where the baseline law holds exactly.
\end{enumerate}
Candidates passing all three gates are optimized via JAX JIT-traced L-BFGS-B and reranked by the Schwarz BIC: $\text{BIC} = N \ln(\text{NMSE}) + k \ln(N)$.

%==================================================================
\section{Experimental Evaluation}
\label{sec:experiments}

\subsection{Synthetic Benchmark: Noise Robustness Study}
We evaluated ADCD across 9 canonical physics scenarios (Relativistic KE, Yukawa Gravity, Anharmonic Spring, Screened Coulomb, Net Radiation, Nonlinear Drag, and three synthetic Mystery functions) at four noise levels ($0\%, 1\%, 5\%, 10\%$). All experiments were repeated across 16 random seeds for statistical reliability.

\begin{figure}[h!]
\centering
\begin{subfigure}[b]{0.48\textwidth}
    \centering
    \includegraphics[width=\textwidth]{fig1_noise_robustness.pdf}
    \caption{Noise robustness compared with PySR}
    \label{fig:noise_robustness}
\end{subfigure}
\hfill
\begin{subfigure}[b]{0.48\textwidth}
    \centering
    \includegraphics[width=\textwidth]{fig4_ablation.pdf}
    \caption{Ablation of Physics Gates}
    \label{fig:ablation}
\end{subfigure}
\caption{(a) Mean structural recovery across 16 seeds: ADCD declines monotonically from $86.8\%$ (0\% noise) to $76.4\%$ (10\% noise); PySR collapses to $11.1\%$ at $5\%$ noise and exhibits non-monotonic behaviour. (b) Component ablation: removing any single physics gate causes a steep accuracy drop, confirming that each gate contributes independently.}
\end{figure}

\textbf{Main results.} At $5\%$ noise, ADCD achieves a mean structural recovery of $\mathbf{77.1\%}$ across 16 seeds (range $69.4\%$--$94.4\%$; bootstrap 95\% CI [$76.7\%$, $84.0\%$]) versus $\mathbf{11.1\%}$ for PySR under matched evaluation criteria, a mean gap of $+65.9$\,pp. A single representative run achieves $88.9\%$ (8/9 scenarios). Under a generous evaluation where PySR is permitted to match simplified sub-expressions, PySR reaches $55.6\%$; ADCD's worst-performing seed ($69.4\%$) still exceeds this threshold.

\textbf{Monotonic noise degradation.} ADCD's accuracy declines smoothly with noise: $86.8\%$, $81.3\%$, $77.1\%$, $76.4\%$ across noise levels $0\%$ to $10\%$. PySR exhibits non-monotonic recovery across the same scenarios, indicating that its stochastic search is not reliably calibrated to noise magnitude.

\textbf{Accuracy versus speed.} ADCD's per-scenario wall-clock time ranges from $1.8$\,s to $14.2$\,s on a single CPU core without parallelism or GPU, which is in the same order of magnitude as PySR under matched settings (100 iterations, 60-second timeout). We do not claim that ADCD is faster than PySR in absolute wall-clock time. However, PySR terminates with \emph{incorrect} functional forms in 8 of 9 scenarios at $5\%$ noise. The relevant metric for scientific discovery is \emph{time-to-correct-structural-recovery}; under this criterion, ADCD's effective throughput is higher because its runs yield correct results more often.

\subsection{Real Physical Scenario Recovery}
To validate recovery on observationally motivated benchmarks, ADCD was applied to four real physical scenarios: (1)~Mercury perihelion relativistic correction, (2)~QED Lamb shift scaling ($\Delta E \propto \alpha^5 m_e c^2 \ln(1/\alpha^2)$), (3)~Muon anomalous magnetic moment Schwinger term ($\alpha/2\pi$), and (4)~Blackbody radiation Wien suppression. ADCD recovered the correct functional class in all four cases (structural match: 4/4). Numerical parameter accuracy varied by scenario; the Muon $g{-}2$ Schwinger prefactor converged to within $3.7\%$ of the theoretical value.

%==================================================================
\section{Real-World Application: SPARC Galaxy Kinematics}
\label{sec:sparc}

To test ADCD on raw, unprocessed observational data, we applied it to the SPARC dataset~\cite{lelli2016sparc}, comprising $3{,}342$ rotation curve measurements across $171$ disc galaxies. Taking Newtonian baryonic acceleration as $f_{\text{base}} = g_{\text{bar}}$ and normalising by the MOND scale $a_0 = 1.2\times10^{-10}$\,m\,s$^{-2}$, ADCD discovered:
\begin{equation}
\nu_{\text{ADCD}}(x) = \theta_0\!\left(\sqrt{1 + \frac{\theta_1}{x}} - 1\right) + 1, \quad x = \frac{g_{\text{bar}}}{a_0}
\label{eq:nu_adcd}
\end{equation}
with bootstrap-mean parameters $\hat{\theta}_0 \approx 2.24$ and $\hat{\theta}_1 \approx 0.27$ (bootstrap 95\% CI: $\theta_0 \in [1.05,\,4.56]$, $\theta_1 \in [0.07,\,0.56]$; single-fit point estimate: $\hat{\theta}_0 = 1.83$, see full technical report). The wide confidence interval on $\theta_0$ reflects genuine parameter degeneracy in the stacked SPARC sample and is reported here transparently.

Algebraically, $\nu_{\text{ADCD}}(x) = (1-\theta_0) + \theta_0\sqrt{1+\theta_1/x}$ belongs to the parametric family of Simple MOND~\cite{famaey2005simple}. The recovered $\hat{\theta}_1 \approx 0.26$ differs from the canonical transition parameter ($c=4.0$), suggesting a delayed Newtonian-to-MOND crossover in the stacked sample. This is a physically interpretable finding that merits further investigation with individual-galaxy fits.

\textbf{Predictive performance.} The framework achieves in-sample NMSE $= 0.373$ and leave-one-out cross-validated NMSE $= 0.386$, outperforming zero-parameter Simple MOND (CV-NMSE $= 0.673$) and the Radial Acceleration Relation (RAR; CV-NMSE $= 0.660$), though statistically indistinguishable from parameter-matched two-parameter refits of these forms (Simple MOND 2-param CV-NMSE $= 0.377$, $z = +0.39$, $p = 0.70$; see full technical report). This improvement over the zero-parameter baselines is achieved without injecting any prior knowledge of MOND physics into the search grammar.

%==================================================================
\section*{AI Assistance and Declarations}

All research concepts, theoretical formulations, experimental design decisions, physical interpretations, and intellectual contributions in this work are entirely the authors' own. AI coding assistants (Google DeepMind's Antigravity agent, Cursor IDE, and LLM-based development environments) were used as pair-programming tools for code implementation, benchmark execution, figure rendering, and manuscript formatting under direct human direction. No AI system independently generated physical hypotheses, selected experimental designs, or produced the reported results.

%==================================================================
\section{Conclusion}
\label{sec:conclusion}

ADCD demonstrates that restricting symbolic search to a physics-anchored correction subspace and pre-screening candidates with analytical gates outperforms unconstrained SR under observational noise. Mean structural recovery across 16 seeds at $5\%$ noise is $77.1\%$ versus $11.1\%$ for PySR under matched computational budget, a $+65.9$\,pp gap achieved on a single CPU core without GPU or parallel compute. ADCD does not claim to be faster than existing SR systems; it claims to be \emph{more accurate} at comparable cost, and to degrade \emph{monotonically} under increasing noise rather than erratically. Functional class recovery on 4/4 real physical scenarios and autonomous discovery of Simple-MOND dynamics from $3{,}342$ raw SPARC observations further support ADCD as a transparent, reproducible, and resource-accessible tool for automated scientific discovery.

\begin{thebibliography}{10}

\bibitem{brunton2016sindy}
S.~L. Brunton, J.~L. Proctor, and J.~N. Kutz.
\newblock Discovering governing equations from data by sparse identification of nonlinear dynamical systems.
\newblock {\em Proceedings of the National Academy of Sciences}, 113(15):3932--3937, 2016.

\bibitem{cranmer2020pysr}
M.~Cranmer.
\newblock PySR: Fast \& parallelized symbolic regression in Python/Julia.
\newblock {\em arXiv preprint arXiv:2009.08295}, 2020.

\bibitem{cranmer2020lagrangian}
M.~Cranmer, A.~Sanchez-Gonzalez, P.~Battaglia, R.~Xu, K.~Cranmer, D.~Spergel, and S.~Ho.
\newblock Lagrangian neural networks.
\newblock {\em ICLR Workshop on Computer Vision and Hardware Acceleration}, 2020.

\bibitem{famaey2005simple}
B.~Famaey and J.~Binney.
\newblock Modified Newtonian dynamics in the Milky Way.
\newblock {\em Monthly Notices of the Royal Astronomical Society}, 363(2):603--608, 2005.

\bibitem{greydanus2019hamiltonian}
S.~Greydanus, M.~Dzamba, and J.~Yosinski.
\newblock Hamiltonian neural networks.
\newblock {\em Advances in Neural Information Processing Systems}, 32, 2019.

\bibitem{lacava2021srbench}
W.~La~Cava, P.~Orzechowski, B.~Burlacu, F.~de~Fran\c{c}a, M.~Virgolin, Y.~Jin, M.~Korn, and J.~Moore.
\newblock Contemporary symbolic regression methods and their performance on synthetic problems.
\newblock {\em Neural Networks}, 143:408--422, 2021.

\bibitem{lelli2016sparc}
F.~Lelli, S.~S. McGaugh, and J.~M. Schombert.
\newblock SPARC: Mass models for 175 disk galaxies with Spitzer photometry and accurate rotation curves.
\newblock {\em The Astronomical Journal}, 152(6):157, 2016.

\bibitem{petersen2021deep}
B.~K. Petersen, M.~L. Larma, T.~N. Mundhenk, C.~P. Santiago, S.~K. Kim, and J.~T. Kim.
\newblock Deep symbolic regression: Recovering mathematical expressions from data via risk-seeking policy gradients.
\newblock {\em International Conference on Learning Representations (ICLR)}, 2021.

\bibitem{raissi2019physics}
M.~Raissi, P.~Perdikaris, and G.~E. Karniadakis.
\newblock Physics-informed neural networks: A deep learning framework for solving forward and inverse problems involving nonlinear partial differential equations.
\newblock {\em Journal of Computational Physics}, 378:686--707, 2019.

\bibitem{rissanen1978mdl}
J.~Rissanen.
\newblock Modeling by shortest data description.
\newblock {\em Automatica}, 14(5):465--471, 1978.

\bibitem{schmidt2009distilling}
M.~Schmidt and H.~Lipson.
\newblock Distilling free-form natural laws from experimental data.
\newblock {\em Science}, 324(5923):81--85, 2009.

\bibitem{tenachi2023physo}
W.~Tenachi, R.~Ibata, and F.~I. Diakogiannis.
\newblock Deep symbolic regression for physics guided by units constraints: Toward the automated discovery of physical laws.
\newblock {\em The Astrophysical Journal}, 959(2):99, 2023.

\bibitem{udrescu2020ai}
S.~M. Udrescu and M.~Tegmark.
\newblock AI Feynman: A physics-inspired method for symbolic regression.
\newblock {\em Science Advances}, 6(16):eaay2631, 2020.

\bibitem{weatheritt2016turbulence}
J.~Weatheritt and R.~Sandberg.
\newblock A novel evolutionary algorithm applied to algebraic Reynolds stress model development.
\newblock {\em Journal of Computational Physics}, 325:22--37, 2016.

\bibitem{rackauckas2020universal}
C.~Rackauckas, Y.~Ma, J.~Martensen, C.~Warner, K.~Zubov, R.~Supekar, D.~Skinner, A.~Ramadhan, and A.~Edelman.
\newblock Universal Differential Equations for Scientific Machine Learning.
\newblock {\em arXiv preprint arXiv:2001.04385}, 2020.

\end{thebibliography}

\end{document}
